\documentclass[11pt]{article}
\usepackage[margin=1in]{geometry}
\usepackage{amsmath,amssymb,amsthm,mathtools}
\usepackage{booktabs,array,enumitem}
\usepackage[T1]{fontenc}
\usepackage{lmodern,microtype}
\usepackage[colorlinks=true,linkcolor=blue,urlcolor=blue,citecolor=blue]{hyperref}
\allowdisplaybreaks
\setlist{itemsep=3pt,topsep=5pt}

\newtheorem{theorem}{Theorem}[section]
\newtheorem{proposition}[theorem]{Proposition}
\newtheorem{lemma}[theorem]{Lemma}
\newtheorem{corollary}[theorem]{Corollary}
\theoremstyle{remark}
\newtheorem{remark}[theorem]{Remark}
\newtheorem{example}[theorem]{Example}

\newcommand{\ii}{\mathrm{i}}
\newcommand{\bbone}{\boldsymbol{1}}
\newcommand{\btheta}{\boldsymbol{\theta}}
\newcommand{\bgamma}{\boldsymbol{\gamma}}
\newcommand{\bk}{\boldsymbol{k}}
\newcommand{\bh}{\boldsymbol{h}}
\newcommand{\bx}{\boldsymbol{x}}
\newcommand{\by}{\boldsymbol{y}}
\newcommand{\be}{\boldsymbol{e}}
\newcommand{\bDelta}{\boldsymbol{\Delta}}
\newcommand{\cH}{\mathcal{H}}
\newcommand{\cD}{\mathcal{D}}
\newcommand{\rtil}{\widetilde{r}}

\title{Optimal Block Weights for Binary Lattice Sequence Prefixes\\[2pt]
\large Path I: one mass per dyadic block, equal point weights within each block}
\author{Megan Millet}
\date{November 11, 2026}

\begin{document}
\maketitle

\begin{abstract}
When integration with a lattice sequence stops at a sample size $n$
that is not a power of two, one can discard the surplus points and keep
the largest dyadic prefix, or weight all $n$ completed values. Within the
weighted compound-rule class of Hickernell, Kritzer, Kuo and Nuyens
(HKKN12), we assign one mass to each dyadic block of the binary
expansion of $n$ and share it equally among that block's points. For
base $2$ and the periodic Korobov Hilbert space, the Gram matrix of the
block error representers has a simple structure: each diagonal entry is a
full-lattice squared error $d_m$, and every off-diagonal entry in a row
equals $2d_{m+1}-d_m$. Consequently the optimal block masses follow from
an $O(\log n)$ scalar recursion once the values $d_0,\dots,d_J$ are cached.
The optimal masses are strictly positive. The squared worst-case error has
an exact expression that is strictly smaller than that of dyadic truncation
at every non-dyadic $n$. We give exact examples and a cancellation-reducing
form of the recursion. We also note two limits: the block optimum cannot
remove the full $(n/2^{\lfloor\log_2 n\rfloor})^\alpha$ loss with a fixed
constant, and the analogous block family gives no improvement over
truncation in the Banach-space figure of merit $P_\alpha$. These results
specialize the HKKN12 framework; novelty and floating-point stability are
not yet established.
\end{abstract}

%======================================================================
\section{Motivation and scope}\label{sec:motivation}
%======================================================================

The referee reports on \emph{The Quality of Lattice Sequences for
Arbitrary Sample Size}~\cite{MDH26} raise a substantive point: an error
rate valid for all $n$ does not by itself justify using all $n$ function
values. Put
\[
 M=\lfloor\log_2 n\rfloor,\qquad N=2^M .
\]
Keeping only the first $N$ values (\emph{truncation}) already attains the
dyadic error bound while discarding fewer than half of the evaluations,
since $n-N<N$. Unrestricted point-by-point optimal weights can
reproduce truncation by putting zero weight on the surplus points, so their
error is trivially no worse than truncation and is non-increasing in $n$.

The question addressed here is therefore:
\begin{quote}
How cheaply can all completed values be used, by how much does doing so
reduce the worst-case error relative to truncation, and can the resulting
weights be guaranteed positive?
\end{quote}
This document develops the \emph{block-constant} answer (Path~I).
The complementary path, which optimizes all $n$ point weights, is not
treated here except for one comparison example
(Section~\ref{sec:fullcompare}).

\paragraph{Relation to HKKN12.}
Sections~3.1, 3.3, 4.2 and 4.3 of~\cite{HKKN12} already provide the
optimal compound weights $A^{-1}\bbone/(\bbone^\top A^{-1}\bbone)$, the
power weights $w_\ell\propto(b^\ell)^a$, the base-$b$ block partition of a
lattice sequence prefix, and optimal compound lattice rules.
Proposition~2 of~\cite{HKKN12} (pp.~174--176) reduces each block
covariance to a single kernel average over the larger block. Neither the
partition $13=8+4+1$ nor the small quadratic program is new. The possible
contributions here are:
\begin{enumerate}[label=(\roman*)]
\item the binary covariance identity of Proposition~\ref{prop:cov}, which
  expresses every entry of $A$ through the cached full-lattice errors
  $d_m$;
\item the $O(L)$ scalar recursion and strict positivity of the optimal
  masses (Theorem~\ref{thm:recursion});
\item the exact gain over truncation (Corollary~\ref{cor:gain});
\item the precomputation and streaming cost analysis
  (Section~\ref{sec:cost}).
\end{enumerate}
No novelty claim has been verified.

\paragraph{Scope.}
The general-base block class is recalled in Section~\ref{sec:general}.
All structural results and proofs concern base $b=2$, completed sequence
prefixes, and the periodic Korobov Hilbert space with $\alpha>1$.
They are not asserted for other bases, arbitrary subsets of points, or the
Banach-space figure of merit $P_\alpha$. Section~\ref{sec:banach} shows that
the Banach analogue genuinely differs.

%======================================================================
\section{Hilbert-space setting and binary lattice sequences}\label{sec:setup}
%======================================================================

\subsection{The periodic Korobov space}
Let $\alpha>1$ and $\gamma_j>0$. The periodic Korobov Hilbert space
$\cH=\mathcal F_{\alpha,2}$ of~\cite{MDH26} has the shift-invariant
reproducing kernel
\begin{equation}\label{eq:kernel}
 K(\bx,\by)=\kappa(\bx-\by),\qquad
 \kappa(\bx)=1+\sum_{\bk\ne\boldsymbol 0}\rho_{\bk}\,
   e^{2\pi\ii\bk^\top\bx},\qquad
 \rho_{\bk}=\rtil(\bk,\bgamma)^{-2\alpha}>0,
\end{equation}
where $\rtil(\bk,\bgamma)=\prod_{j:k_j\ne0}|k_j|/\gamma_j$ (the empty
product is $1$) and $\ii^2=-1$. The series converges absolutely, $\kappa$
is real and even, and $\int_{[0,1)^s}K(\bx,\by)\,d\bx=1$.
For integer $\alpha$ the kernel has the closed product form
\begin{equation}\label{eq:bernoulli}
 \kappa(\bx)=\prod_{j=1}^s\left[1+
 \frac{(-1)^{\alpha+1}(2\pi\gamma_j)^{2\alpha}}{(2\alpha)!}
 B_{2\alpha}(\{x_j\})\right],
\end{equation}
where $B_{2\alpha}$ is the Bernoulli polynomial of degree $2\alpha$ and
$\{\cdot\}$ denotes the fractional part.

For a rule $Q(f)=\sum_i w_if(\bx_i)$, let
$e(Q)=\sup_{\|f\|_\cH\le1}|Q(f)-I(f)|$. If $\sum_iw_i=1$, then
\begin{equation}\label{eq:fourier}
 e(Q)^2=\sum_{\bk\ne\boldsymbol 0}\rho_{\bk}
 \Bigl|\sum_iw_i\,e^{2\pi\ii\bk^\top\bx_i}\Bigr|^2
 = w^\top G w-1,
\end{equation}
where $G=(\kappa(\bx_i-\bx_j))_{i,j}$ is the pointwise Gram matrix.
A common shift $\bDelta$ of all nodes does not change~\eqref{eq:fourier}.

\begin{remark}[Normalization]
For weights that do not sum to one, the correct identity is
$e(Q)^2=1-2\bbone^\top w+w^\top Gw$. The simplification to
$w^\top Gw-1$ requires $\bbone^\top w=1$. All weights below are normalized.
\end{remark}

\subsection{Binary lattice sequences and dyadic errors}
Let $\phi_2$ denote the van der Corput radical inverse in base $2$, and let
\[
 \bx_i=\{\bh\,\phi_2(i)+\bDelta\},\qquad i=0,1,2,\dots,
\]
be a rank-1 lattice sequence with at least one odd generator coordinate,
say $h_1$. All prefixes considered lie within a finite dyadic horizon
$2^J$ (fixed in Section~\ref{sec:cost}); on that horizon we use one
coherent integer representative of $\bh$ modulo $2^J$.
(An existence theorem for a good infinite generator does not certify any
particular software default.) Define
\[
 H_m=\{\bh j/2^m \bmod 1: 0\le j<2^m\},\qquad
 \cD_m=\{\bk\ne\boldsymbol 0:\bk^\top\bh\equiv0 \pmod{2^m}\}.
\]
Since $h_1$ is odd, $|H_m|=2^m$, and $H_m\subset H_{m+1}$.
The squared worst-case error of the full $2^m$-point equal-weight rule is
\begin{equation}\label{eq:dm}
 d_m:=e(Q_{2^m})^2
 =\frac1{2^m}\sum_{j=0}^{2^m-1}\kappa\!\left(\frac{\bh j}{2^m}\right)-1
 =\sum_{\bk\in\cD_m}\rho_{\bk}.
\end{equation}
Because $\cD_{m+1}\subset\cD_m$ and $\bk=2^m\be_1\in\cD_m\setminus\cD_{m+1}$
(as $h_1$ is odd), we have the strict ordering
\begin{equation}\label{eq:dorder}
 d_m>d_{m+1}>0,\qquad
 d_m-d_{m+1}=\sum_{\bk\in\cD_m\setminus\cD_{m+1}}\rho_{\bk}>0 .
\end{equation}

\paragraph{Nodes are distinct.}
The first $2^J$ values $\phi_2(i)$ are the distinct multiples of $2^{-J}$,
and because $h_1$ is odd their images $\{h_1\phi_2(i)\}$ are distinct.
Hence the first $n\le 2^J$ nodes are distinct, and $G$ is positive definite
because every $\rho_{\bk}>0$.

%======================================================================
\section{The block-constant weight class}\label{sec:general}
%======================================================================

\subsection{General base (the HKKN12 class)}
For an integer base $b\ge2$, put $M_b=\lfloor\log_bn\rfloor$, write
$n=\sum_{m=0}^{M_b}n_mb^m$ with $0\le n_m<b$, and set
$a_m=\sum_{\ell>m}n_\ell b^\ell$. Level $m$ contains $n_m$ separate blocks
of $b^m$ consecutive points, with normalized block rules
\[
 Q_{m,q}(f)=\frac1{b^m}\sum_{i=a_m+qb^m}^{a_m+(q+1)b^m-1}f(\bx_i),
 \qquad 0\le q<n_m.
\]
Assigning a total mass $\theta_{m,q}$ to each block gives
\[
 Q_{\btheta}=\sum_{m=0}^{M_b}\sum_{q=0}^{n_m-1}\theta_{m,q}Q_{m,q},
 \qquad \sum_{m,q}\theta_{m,q}=1 .
\]
Every point in block $(m,q)$ receives weight $\theta_{m,q}/b^m$.
This is the compound-rule class of~\cite{HKKN12}, with
$R=\sum_mn_m\le(b-1)(M_b+1)$ unknowns. If $A$ is the Gram matrix of the
block error functionals $Q_{m,q}-I$ in the dual space $\cH^*$, then
$e(Q_{\btheta})^2=\btheta^\top A\btheta$. The block optimum is
$\btheta=u/(\bbone^\top u)$, where $Au=\bbone$.

\subsection{Binary blocks}
From now on $b=2$. Write the nonzero binary digits of $n$ as
\[
 n=N_1+\cdots+N_L,\qquad N_r=2^{m_r},\qquad m_1>\cdots>m_L\ge0,
 \qquad a_r=\sum_{q<r}N_q .
\]
Thus $L\le M+1$, $m_1=M$, block $r$ is $B_r=\{a_r,\dots,a_r+N_r-1\}$, and
$Q_r(f)=N_r^{-1}\sum_{i\in B_r}f(\bx_i)$. The block-constant rule is
\begin{equation}\label{eq:compound}
 Q_{\btheta}(f)=\sum_{r=1}^L\theta_rQ_r(f),\qquad
 \sum_{r=1}^L\theta_r=1,\qquad
 w_i=\frac{\theta_r}{N_r}\quad(i\in B_r).
\end{equation}
The $\theta_r$ are \emph{block masses}, not individual sample weights.
Familiar members of this class are:
\begin{center}
\begin{tabular}{@{}ll@{}}
\toprule
Rule & Block masses $\theta_r$\\
\midrule
Equal point weights & $N_r/n$\\
HKKN12 power weights, $a>0$ & $N_r^a\big/\sum_qN_q^a$\\
Truncation to $2^M$ points & $(1,0,\dots,0)$\\
\bottomrule
\end{tabular}
\end{center}

\subsection{The binary block covariance}

\begin{proposition}[Binary block covariance]\label{prop:cov}
Let $A_{rq}=\langle Q_r-I,\,Q_q-I\rangle_{\cH^*}$, with blocks in
decreasing size order. Then
\begin{equation}\label{eq:cov}
 A_{rr}=d_{m_r},\qquad
 A_{rq}=A_{qr}=c_{m_r}:=2d_{m_r+1}-d_{m_r}\quad(r<q).
\end{equation}
In particular, all off-diagonal entries in row $r$ to the right of the
diagonal are equal, and $e(Q_{\btheta})^2=\btheta^\top A\btheta$.
\end{proposition}

\begin{proof}
Since $\int K(\bx,\by)\,d\bx=1$,
\[
 A_{rq}=\frac1{N_rN_q}\sum_{i\in B_r}\sum_{j\in B_q}
 \kappa(\bx_i-\bx_j)-1 .
\]
By the block property of the van der Corput sequence, block $B_r$ is
$\{\bx_{a_r}+\by:\by\in H_{m_r}\}$ (mod $1$), a full coset of $H_{m_r}$.
The diagonal entry is therefore $d_{m_r}$, by shift invariance.

Let $r<q$. Because $H_{m_q}\subset H_{m_r}$, block $B_q$ lies in a single
coset of $H_{m_r}$. The function
$\by\mapsto N_r^{-1}\sum_{i\in B_r}\kappa(\bx_i-\by)$ is constant on cosets
of $H_{m_r}$. So
\[
 A_{rq}=\frac1{N_r}\sum_{j=0}^{N_r-1}
 \kappa\!\left(\frac{\bh j}{N_r}+\bh\bigl(\phi_2(a_q)-\phi_2(a_r)\bigr)\right)-1 .
\]
Now $a_r$ is a multiple of $2N_r$, and
$a_q=a_r+N_r+\sum_{r<p<q}N_p$ with $\sum_{r<p<q}N_p<N_r$. In binary, $a_q$
and $a_r$ therefore agree in all digits above position $m_r$. The digit at
position $m_r$ changes from $0$ to $1$ and contributes $2^{-(m_r+1)}$ to the
radical inverse. Digits of $a_q$ below position $m_r$ contribute multiples
of $2^{-m_r}$. Hence
\[
 \phi_2(a_q)-\phi_2(a_r)\equiv\frac1{2N_r}\pmod{\frac1{N_r}}.
\]
Since $\kappa$ is $1$-periodic and $\bh$ is integral on the horizon, the
residual multiple of $1/N_r$ only permutes the index $j$. Thus
\[
 A_{rq}=\frac1{N_r}\sum_{j=0}^{N_r-1}
 \kappa\!\left(\bh\,\frac{j+1/2}{N_r}\right)-1 .
\]
The lattice $H_{m_r+1}$ is the disjoint union of $H_{m_r}$ and
$H_{m_r}+\bh/(2N_r)$. Averaging $\kappa-1$ over $H_{m_r+1}$ gives
$d_{m_r+1}=\tfrac12\bigl(d_{m_r}+A_{rq}\bigr)$, which is~\eqref{eq:cov}.
\end{proof}

\begin{remark}
Proposition~\ref{prop:cov} specializes~\cite[Proposition~2]{HKKN12} and
should be cited as such. The extra binary information is that the
cross-covariance does not depend on $q$ and is determined by the cached
dyadic errors. No Toeplitz structure of the pointwise Gram matrix (in
van der Corput order) is assumed or needed.
\end{remark}

%======================================================================
\section{Optimal block masses}\label{sec:optimal}
%======================================================================

The general formula $A^{-1}\bbone/(\bbone^\top A^{-1}\bbone)$ applies.
Equation~\eqref{eq:cov} allows the same optimum to be computed without
forming $A$ or solving a linear system.

\begin{theorem}[Optimal binary block rule]\label{thm:recursion}
The minimizer of $\btheta^\top A\btheta$ over real $\btheta$ with
$\bbone^\top\btheta=1$ is unique, and every component is strictly
positive. It is computed as follows. Set $E_L=d_{m_L}$. For
$r=L-1,\dots,1$, compute
\begin{align}
 D_r&=d_{m_r}+E_{r+1}-2c_{m_r},\label{eq:denom}\\
 t_r&=\frac{E_{r+1}-c_{m_r}}{D_r},\label{eq:tr}\\
 E_r&=\frac{d_{m_r}E_{r+1}-c_{m_r}^{\,2}}{D_r}.\label{eq:Er}
\end{align}
Recover the masses in a forward pass:
\[
 p\leftarrow1;\quad
 \theta_r=p\,t_r,\ \ p\leftarrow p(1-t_r)\quad(r=1,\dots,L-1);\quad
 \theta_L=p .
\]
If $L=1$, then $\theta_1=1$. The optimal squared worst-case error is $E_1$.
More generally, $E_r$ is the optimal squared error among normalized rules
supported on blocks $r,\dots,L$.
\end{theorem}

\begin{proof}
\emph{Step 1: a lower bound for tails.}
Let $R$ be any normalized block-constant rule (signed weights allowed)
supported on blocks $r+1,\dots,L$. By the proof of
Proposition~\ref{prop:cov}, all its nodes lie in one coset of $H_{m_r}$.
For $\bk\in\cD_{m_r}$, the character $e^{2\pi\ii\bk^\top\bx}$ is constant on
each coset of $H_{m_r}$. Hence
$|R(e^{2\pi\ii\bk^\top\cdot})|=|\sum_iw_i|=1$, and by~\eqref{eq:fourier}
\begin{equation}\label{eq:tail}
 e(R)^2\ \ge\ \sum_{\bk\in\cD_{m_r}}\rho_{\bk}=d_{m_r}.
\end{equation}
By~\eqref{eq:cov}, $\langle Q_r-I,R-I\rangle=c_{m_r}$, whatever the tail
masses are.

\emph{Step 2: one-dimensional reduction.}
Any normalized rule on blocks $r,\dots,L$ either has tail mass $1-t\ne0$,
in which case it equals $tQ_r+(1-t)R$ with $R$ a normalized tail rule, or
has tail mass zero. In the first case,
\[
 e^2=t^2d_{m_r}+2t(1-t)c_{m_r}+(1-t)^2e(R)^2 ,
\]
which, for fixed $t$, is minimized by taking $e(R)^2=E_{r+1}$ (the inductive
hypothesis). Write $v=E_{r+1}$. By~\eqref{eq:tail} and~\eqref{eq:dorder},
$v\ge d_{m_r}>c_{m_r}$. So the quadratic coefficient
$D_r=d_{m_r}+v-2c_{m_r}$ is positive, and the minimizer
$t_r=(v-c_{m_r})/D_r$ satisfies $0<t_r<1$, because $v-c_{m_r}>0$ and
$D_r-(v-c_{m_r})=d_{m_r}-c_{m_r}>0$. Substituting gives~\eqref{eq:Er}.

In the second case the rule is $Q_r+Z$, where $Z=\sum_{p>r}\theta_pQ_p$ and
$\sum_{p>r}\theta_p=0$. The cross term is
$\langle Q_r-I,Z\rangle=c_{m_r}\sum_{p>r}\theta_p=0$, so $e^2=d_{m_r}+\|Z\|^2\ge
d_{m_r}$. This is no better than $t=1$, which is strictly worse than the
interior optimum, whose value is
$E_r=d_{m_r}-(d_{m_r}-c_{m_r})^2/D_r<d_{m_r}$.

\emph{Step 3: induction and positivity.}
Backward induction from $E_L=d_{m_L}$ proves the recursion. Each $\theta_r$
is a product of factors $t_r\in(0,1)$ and $1-t_q\in(0,1)$, so every mass is
positive.

\emph{Uniqueness.}
$A$ is positive definite. Suppose a nonzero $u$ gave
$\sum_ru_r(Q_r-I)=0$ on $\cH$. Then the finite signed atomic measure
$\sum_ru_rN_r^{-1}\sum_{i\in B_r}\delta_{\bx_i}$ and
$(\sum_ru_r)$ times Lebesgue measure would have identical Fourier
coefficients, because every $\rho_{\bk}>0$. They would therefore be equal,
which is impossible for a nonzero atomic measure on distinct nodes.
Hence the strictly convex problem has a unique minimizer.
\end{proof}

\begin{corollary}[Exact gain over truncation]\label{cor:gain}
Let $n$ be non-dyadic, so $L\ge2$ and $m_1=M$. Then
\begin{equation}\label{eq:gain}
 e(Q_{\rm block})^2=E_1
 =d_M-\frac{4(d_M-d_{M+1})^2}{(E_2-d_M)+4(d_M-d_{M+1})}
 <d_M=e(Q_{\rm trunc})^2 .
\end{equation}
All $n$ point weights are positive. Moreover,
\[
 e(Q_{\rm block})\le e(Q_{\rm HKKN}(a))\ \text{for every }a>0,\qquad
 e(Q_{\rm block})\le e(Q_{\rm equal}),
\]
and the unrestricted point-weight optimum satisfies
$e(Q_{\rm full\ opt})\le e(Q_{\rm block})$.
If $e(Q_{2^m})\le C\,2^{-\beta m}$ for all relevant $m$, then
$e(Q_{\rm block})\le C\,2^\beta n^{-\beta}$.
\end{corollary}

\begin{proof}
The denominator equals $D_1=d_M+E_2-2c_M$, because
$d_M-c_M=2(d_M-d_{M+1})$. The identity
$E_1=d_M-(d_M-c_M)^2/D_1$ is the minimum value of the quadratic in the
proof of Theorem~\ref{thm:recursion}. The inequalities follow because each
compared rule is feasible for the corresponding minimization. For the rate,
use $e(Q_{\rm block})\le e(Q_{\rm trunc})\le C2^{-\beta M}$ and $2^M>n/2$.
\end{proof}

\begin{remark}[What the corollary does and does not say]
The rate is already available from truncation and from~\cite{HKKN12}.
The substantive statements are positive use of every point, a
\emph{quantified} strict gain, and an inexpensive computation.
Worst-case inequalities do not order the actual errors for an individual
integrand. They also do not imply that unrestricted point-wise optimal
weights are positive.
\end{remark}

\begin{remark}[Connection to the research questions]
In the notation of the earlier draft on optimal block weights,
Theorem~\ref{thm:recursion} answers (Q1): in base $2$ the optimal block
weights have a closed recursive characterization, computable in $O(L)$
operations. Equation~\eqref{eq:gain} is an exact error value that keeps
all inter-block cancellation, which bears on (Q3). Since the gain depends
only on $d_M$, $d_{M+1}$ and the tail error, it also suggests a criterion
relevant to generator design (Q5): the differences $d_m-d_{m+1}$ should not
be too small relative to $d_m$. This is a conjecture, not a result.
\end{remark}

%======================================================================
\section{Precomputation, cost and numerical stability}\label{sec:cost}
%======================================================================

\subsection{Recursion versus dense solves}
The backward recursion and forward pass cost $O(L)=O(\log n)$ arithmetic
and storage, given the cached $d_m$. They return the $L$ block masses,
which encode all point weights through $w_i=\theta_r/N_r$. Writing out the
full $n$-vector of point weights would itself cost $O(n)$.

For comparison, solving $Au=\bbone$ by dense Cholesky costs $O(L^3)$
arithmetic and $O(L^2)$ storage, excluding assembly. A dense solve for the
unrestricted $n$-point optimum costs $O(sn^2+n^3)$ arithmetic and $O(n^2)$
storage, and generally produces a different rule. These figures are
baselines, not lower bounds.

\subsection{One reusable cache}
For a maximum prefix length $n_{\max}\ge1$, set
$J=\lceil\log_2n_{\max}\rceil$ and $H=2^J$. The cache $d_0,\dots,d_J$
supports every $1\le n\le n_{\max}$. A non-dyadic $n$ needs $d_{m_r+1}$
with $m_r+1\le\lfloor\log_2n\rfloor+1\le J$, and a dyadic $n$ needs only
$d_{\log_2n}$.

The lattices are nested, so the kernel needs evaluating only at newly
added nodes:
\[
 S_0=\kappa(\boldsymbol 0),\qquad
 S_{m+1}=S_m+\sum_{j=0}^{2^m-1}\kappa\!\left(\bh\,\frac{2j+1}{2^{m+1}}\right),
 \qquad d_m=2^{-m}S_m-1 .
\]
This requires exactly $1+\sum_{m=0}^{J-1}2^m=H$ kernel evaluations, which
is $O(sH)=O(sn_{\max})$ arithmetic when one kernel evaluation costs $O(s)$.
Streaming the sums uses $O(s+J)$ working storage, and neither the nodes nor
a kernel matrix need be stored. The cache depends on the kernel and
generator only. It does not depend on the integrand, the shift or $n$, and
it needs no integrand evaluations beyond the completed prefix. If
$n_{\max}$ is not known in advance, the cache can be extended one level at
a time by the same nested update.

\subsection{Streaming estimates}
Keep one function-value sum per active binary block. When a new value
arrives, append a block of size one and merge adjacent equal-size blocks,
as in a binary counter. Through $n_{\max}$ samples the merges cost
$O(n_{\max})$ arithmetic in total and use $O(\log n_{\max})$ storage.
Each requested estimate costs $O(L)$ once the masses are known.
Reporting after every sample therefore costs $O(n_{\max}\log n_{\max})$
overall, in addition to the reusable $O(sn_{\max})$ precomputation and the
integrand evaluations.

The data must form a completed sequence prefix. Lost values or an arbitrary
asynchronous subset of points is a different problem.

\subsection{Cancellation and numerical validation}
The operation counts above say nothing about numerical stability. Three
cancellation risks are known.

\paragraph{(a) Subtracting the unit baseline.}
The formula $d_m=2^{-m}S_m-1$ extracts a small residual from an average
close to one. An absolute error $\eta$ in the average becomes a relative
error of about $|\eta|/d_m$. For example, with $s=1$, $h=1$, $\alpha=2$ and
$\gamma=1$, $d_{14}=2\zeta(4)\,2^{-56}\approx3.0\times10^{-17}$, so even the
exact average $1+d_{14}$ rounds to $1$ in binary64. Compensated summation
cannot recover a residual already lost in rounding. (In our random
multidimensional tests in binary64, entries of $A$ formed as ``average
minus one'' agreed with~\eqref{eq:cov} only to a relative accuracy of about
$10^{-5}$. In 60-digit arithmetic the agreement was about $10^{-57}$.)

A first remedy~\cite[Section~4.4]{JH19} is to evaluate the centered kernel
directly. If $\kappa(\bx)=\prod_{j=1}^s[1+a_j(x_j)]$, compute each $a_j$
directly (never as a rounded factor minus one) and accumulate
\begin{equation}\label{eq:centered}
 r_1=a_1,\qquad r_j=r_{j-1}(1+a_j)+a_j\ \ (2\le j\le s),\qquad
 r_s=\kappa-1,
\end{equation}
then form $d_m=2^{-m}\sum_{i<2^m}r_s(\{\bh i/2^m\})$. This costs the same
$O(sH)$. The residuals can have either sign, so cancellation inside the
lattice average still requires validation.

\paragraph{(b) Subtractions in the recursion.}
The determinant-type subtraction in~\eqref{eq:Er} and the forward
subtraction $1-t_r$ can be avoided. Set
\[
 d=d_{m_r},\quad d_+=d_{m_r+1},\quad g=d-d_+>0,\quad
 \delta=E_{r+1}-d\ge0,\quad D=\delta+4g .
\]
Then, exactly,
\begin{equation}\label{eq:stable}
 t_r=\frac{\delta+2g}{D},\qquad
 q_r:=1-t_r=\frac{2g}{D},\qquad
 E_r=d_++g\,\frac{\delta}{D}=d-\frac{4g^2}{D}.
\end{equation}
Compute $q_r$ directly and update $p\leftarrow p\,q_r$ in the forward pass.
All quantities in~\eqref{eq:stable} are nonnegative combinations once $g$
and $\delta$ are known.

\paragraph{(c) Accurate differences.}
Formula~\eqref{eq:stable} moves the difficulty into the inputs
$g=d-d_+$ and $\delta=E_{r+1}-d$. These need sufficiently accurate cached
values. One option is the positive spectral representation~\eqref{eq:dorder}
of $d_m-d_{m+1}$ with a controlled truncation error; another is
higher-precision precomputation. Either may raise the precomputation cost.
Very small masses may also underflow.

The equivalent formulas are not a stability theorem. A stable
implementation and precision-convergence checks are required before
reporting floating-point weights or error values as reliable.

%======================================================================
\section{Examples and limitations}\label{sec:examples}
%======================================================================

Throughout this section $s=1$, $h=1$ and Hilbert smoothness $\alpha=2$,
so that~\eqref{eq:dm} gives
\[
 d_m=C\,2^{-4m},\qquad C=2\gamma^4\zeta(4)=\frac{\gamma^4\pi^4}{45},
 \qquad c_m=2d_{m+1}-d_m=-\tfrac78\,C\,2^{-4m}.
\]
Since the weights are normalized, the optimal weights do not depend on $C$.
All values below were verified in exact rational arithmetic.

\subsection{A thirteen-point example}
For $13=8+4+1$ ($m=3,2,0$), Theorem~\ref{thm:recursion} gives
$t_2=137/138$, $t_1=613/682$ and
\[
 (\theta_1,\theta_2,\theta_3)
 =\left(\frac{613}{682},\ \frac{137}{1364},\ \frac1{1364}\right).
\]
Each of the first eight points has weight $613/5456$, each of the next four
has weight $137/5456$, and the last point has weight $4/5456$. Exactly,
\[
 e(Q_{\rm block})^2=C\,\frac{4421}{22347776},\qquad
 \frac{e(Q_{\rm block})}{e(Q_{\rm trunc})}=\sqrt{\frac{4421}{5456}}
 \approx0.9002 .
\]
This is a sanity check, not high-dimensional evidence. The weights on
surplus points can be very small; positive use of a point does not by
itself mean a large practical benefit.

\subsection{No universal factor-free bound}
At $n=3=2+1$ the optimal masses are $(9/10,1/10)$, the point weights are
$(9/20,9/20,1/10)$, and $e(Q_{\rm block})^2=13C/256$. The dyadic law is
$e(Q_{2^m})=\sqrt C\,(2^m)^{-2}$. Keeping the same constant at $n=3$ would
require $e^2\le C/81$, but $13/256>1/81$. Thus~\eqref{eq:gain} improves the
truncated Hilbert error without removing the whole
$(n/2^M)^{\alpha}$ loss with a fixed constant.

\subsection{Comparison with unrestricted point weights}\label{sec:fullcompare}
The block class is contained in the class of all normalized point weights,
so $e_{\rm full}\le e_{\rm block}$. At $n=5=4+1$ the first five points are
$(0,1/2,1/4,3/4,1/8)$ before a common shift, and
\[
 w_{\rm block}=\tfrac1{552}(137,137,137,137,4),\qquad
 w_{\rm full}=\tfrac1{544}(117,139,117,139,32),
\]
\[
 e_{\rm block}^2=C\,\frac{363}{94208},\qquad
 e_{\rm full}^2=C\,\frac{121}{34816},\qquad
 \frac{e_{\rm full}^2}{e_{\rm block}^2}=\frac{46}{51}.
\]
For further comparison, $e_{\rm full}^2/e_{\rm block}^2$ equals $1$ at
$n=3$ and $n=6$, $3409/5829\approx0.585$ at $n=7$, and
$1219633/1560613\approx0.781$ at $n=13$. These values show that the gap
in (Q2) can be substantial at some $n$ and zero at others. They do not
establish a typical gain, and they say nothing general about the signs of
the unrestricted weights (all were positive at these small $n$).

\subsection{A Banach-space obstruction}\label{sec:banach}
The Hilbert-space improvement does not transfer to the Banach figure of
merit $P_\alpha$ of~\cite{MDH26}. In one dimension, with
$\rtil(k,\gamma)=|k|/\gamma$ for $k\ne0$, put $C_\alpha=2\gamma^\alpha\zeta(\alpha)$
and $q=2^{-\alpha}$. Then $q<1/2$ because $\alpha>1$.
(If the alternative form $\max(1,|k|/\gamma)$ is used, the same computation
requires $\gamma\le1$.)
For $n=3$, give mass $1-t$ to the first two points and mass $t$ to the
singleton. Splitting the wavenumbers into odd integers, integers $\equiv2
\pmod 4$, and nonzero multiples of $4$ gives
\begin{equation}\label{eq:banach}
 P_\alpha(t)=C_\alpha\bigl[q^2+q(1-q)|1-2t|+(1-q)|t|\bigr].
\end{equation}
Its slope is $-C_\alpha(1-q)(1+2q)<0$ for $t<0$,
$C_\alpha(1-q)(1-2q)>0$ for $0<t<1/2$, and $C_\alpha(1-q)(1+2q)>0$ for
$t>1/2$. So the unique minimizer over real $t$ is $t=0$, which is
truncation. In the Banach setting, block weighting therefore cannot be
guaranteed to improve on truncation at every non-dyadic $n$, although it
always does so in the Hilbert setting. A referee's displayed bound concerns
$P_\alpha$; the results here address only its Hilbert analogue
$P_{\alpha,2}$ and do not sharpen any $P_\alpha$ bound.

%======================================================================
\section{Proposed numerical study}\label{sec:experiments}
%======================================================================

The following is a protocol, not completed work.
\begin{enumerate}
\item \textbf{Methods.} On identical nodes and the same kernel, compare
  equal weights, truncation at $2^{\lfloor\log_2n\rfloor}$, HKKN12 power
  masses with $a=2,3,4$, and the optimal block rule of
  Theorem~\ref{thm:recursion}, all normalized. At small $n$, check the
  recursion against the dense $L\times L$ solve of $Au=\bbone$; they must
  agree. Where feasible, include the dense unrestricted optimum as a
  reference for (Q2).
\item \textbf{Parameters.} Use Hilbert $\alpha=2$ as the primary case,
  add $\alpha=3$, and keep $\alpha=1$ only as context. Use, for example,
  $s=2,6,20$ with two declared coordinate-weight decays. Scan every $n$ up
  to about $2^{12}$, then prespecified larger budgets, including
  $2^M+1$, $2^{M+1}-1$, $13\cdot2^p$ and different binary digit counts.
  Report results for a common generator separately from any tuned
  generators.
\item \textbf{Outputs.} Plot $e_{\rm method}/e_{\rm trunc}$ against $n$
  and against $n/2^{\lfloor\log_2n\rfloor}$, marking powers of two. Report
  ratio distributions and maxima, not only selected curves, and keep $e$
  and $e^2$ distinct. Time precomputation, weight computation and
  application separately, and record memory and precision.
\item \textbf{Integrands.} Use smooth periodic products
  $f(\bx)=\prod_j[1+\eta_j\cos(2\pi x_j)]$ (integral one) and a
  broader-spectrum periodic example with a verified finite norm.
  For actual-error comparisons use paired common random shifts (at least
  $32$), with absolute-error summaries. The Keister integrand may serve as
  an illustration only after its regularity and norm in the theoretical
  space are checked. Guard actual-error ratios against a zero
  denominator.
\item \textbf{Checks.} Verify $\sum_iw_i=1$, positivity,~\eqref{eq:cov},
  and the stationarity residual $A\btheta-E_1\bbone$. Run
  precision-convergence checks on $d_m$, $g$, $\delta$ and the masses,
  and never silently clip negative round-off values. Record generators
  (or an archived file with checksum), horizon, ordering, kernel scaling,
  coordinate weights, seeds, software, precision and scripts.
\end{enumerate}

\paragraph{Existing code.}
Two issues in the current numerical code should be fixed before any
figure is reused. First, \texttt{Programs/discrepancy.m} solves
$u=K^{-1}\bbone$ and evaluates the error without normalizing $u$; the
constrained optimum is $w=u/(\bbone^\top u)$, computed by a linear solve
rather than an explicit inverse. Second, \texttt{GramMat.m} uses a $B_2$
factor with coefficient $\texttt{coordwts}_j/2$; under the $\alpha=1$
convention of~\eqref{eq:bernoulli} the coefficient is $2\pi^2\gamma_j^2$, so
the two agree only when $\texttt{coordwts}_j=4\pi^2\gamma_j^2$. An $\alpha=2$
study needs the $B_4$ factor of~\eqref{eq:bernoulli}.

%======================================================================
\section{Status and open questions}\label{sec:status}
%======================================================================

Proposition~\ref{prop:cov}, Theorem~\ref{thm:recursion} and
Corollary~\ref{cor:gain} are drafts for joint verification. Their
algebraic content and all displayed examples were checked independently
as follows. The examples were reproduced in exact rational arithmetic.
The covariance formula, the recursion (against dense solves), positivity,
and the cancellation-reducing form~\eqref{eq:stable} were checked on
random multidimensional cases with odd $h_1$, $\alpha\in\{1,2\}$ and random
shifts. The covariance formula was also confirmed in 60-digit
arithmetic. These checks verify identities on tested cases; they do not
replace the proofs or establish floating-point stability.

Open items:
\begin{enumerate}
\item Check whether the binary specialization and positivity result
  already appear in the literature. Natural comparisons are~\cite{HKKN12},
  symmetry reductions in kernel quadrature~\cite{KS18}, and positivity
  results for Bayesian quadrature weights~\cite{KKS19}, whose hypotheses
  and normalization must be checked before any conclusion is imported.
\item Establish a stable implementation and its precision requirements.
\item Quantify the empirical gain over truncation and over HKKN12 power
  weights, and the gap to the unrestricted optimum (Q2).
\item Investigate whether finer partitions (Q4) or generators designed for
  the compound rule (Q5) reduce the staircase effect.
\end{enumerate}

\begin{thebibliography}{99}
\bibitem{HKKN12}
F.~J. Hickernell, P.~Kritzer, F.~Y. Kuo and D.~Nuyens,
Weighted compound integration rules with higher order convergence for
all $N$,
\emph{Numerical Algorithms} \textbf{59} (2012), 161--183.
\href{https://doi.org/10.1007/s11075-011-9482-5}{doi:10.1007/s11075-011-9482-5}.

\bibitem{HN03}
F.~J. Hickernell and H.~Niederreiter,
The existence of good extensible rank-1 lattices,
\emph{Journal of Complexity} \textbf{19} (2003), 286--300.

\bibitem{MDH26}
L.~Matiukha, Y.~Ding and F.~J. Hickernell,
The quality of lattice sequences for arbitrary sample size,
manuscript under revision (2026).

\bibitem{JH19}
R.~Jagadeeswaran and F.~J. Hickernell,
Fast automatic Bayesian cubature using lattice sampling,
\emph{Statistics and Computing} \textbf{29} (2019), 1215--1229.
\href{https://doi.org/10.1007/s11222-019-09895-9}{doi:10.1007/s11222-019-09895-9}.

\bibitem{KS18}
T.~Karvonen and S.~S\"arkk\"a,
Fully symmetric kernel quadrature,
\emph{SIAM Journal on Scientific Computing} \textbf{40} (2018),
A697--A720.
\href{https://doi.org/10.1137/17M1121779}{doi:10.1137/17M1121779}.

\bibitem{KKS19}
T.~Karvonen, M.~Kanagawa and S.~S\"arkk\"a,
On the positivity and magnitudes of Bayesian quadrature weights,
\emph{Statistics and Computing} \textbf{29} (2019), 1317--1333.
\href{https://doi.org/10.1007/s11222-019-09901-0}{doi:10.1007/s11222-019-09901-0}.
\end{thebibliography}

\end{document}